\section{Arithmetic estimates for the packets}\label{sec:arithmetic}

This section establishes the three estimates on which the dynamical
argument rests: uniform control of the cusp, a dimension bound for
probability limits, and exponential decay on two-sided diagonal tubes.
The first two follow by counting ordinary lattice vectors.  For the
third, we count an ordinary vector together with a dual vector; their
coordinatewise product lies in the codifferent and can be counted in a
thin slice.

For the volume, unfolding, and ordinary-vector estimates, we adapt
the method of \cite[Sections~2--6]{PrimeDegreePackets} and
\cite[Section~3]{QuarticPackets}.  We reproduce the needed arguments
in the packet and Haar normalizations of this paper.

Throughout this section, $K$ is a totally real field of degree six with
no proper intermediate field, and $\sigma$ is an ordering of its real
embeddings.  Write
\[
 D=|\Disc(K)|,\qquad Q=D^{1/2},\qquad
 \kappa=\mathop{\rm Res}_{s=1}\zeta_K(s),\qquad
 d_K(l)=\#\{\mathfrak b\subset\mathcal O_K:N\mathfrak b=l\}.
\]
All ideals counted here are nonzero.  For a sequence of fields and
embedding orderings, we use the corresponding notation
\[
 D_i=|\Disc(K_i)|,\qquad Q_i=D_i^{1/2},\qquad
 \kappa_i=\mathop{\rm Res}_{s=1}\zeta_{K_i}(s),\qquad
 \mu_i=\mu_{K_i,\sigma_i}.
\]
Unless another dependence is indicated, constants below are absolute
for this fixed degree.  In particular, they are independent of the
field and of the ordering of its embeddings.

\subsection{Volume and unfolding}

We first make the sign convention in the packet definition explicit.
The space $X=\SL_6(\Z)\backslash\SL_6(\R)$ parametrizes ordinary
unimodular lattices: a lattice has a positively oriented basis, and
any two such bases differ by an element of $\SL_6(\Z)$.  Consequently
a diagonal sign matrix acts on lattices even when its determinant is
$-1$.  The resulting lattice again represents a point of $X$ by choosing
a positively oriented basis.  There is no additional orientation
coordinate.

Let $h_K$ and $R_K$ denote the ordinary class number and regulator,
respectively.  Our convention for $R_K$ is the covolume of the absolute
logarithms of the units after deleting the sixth coordinate.  Put
\[
 U=\mathcal O_K^\times,\qquad
 U^+=\{\epsilon\in U:\sigma_j(\epsilon)>0\ (1\le j\le6)\},
 \qquad s=[U:U^+].
\]
For a nonnegative Borel function $f$ on $\R^6$, define its lattice sum by
\[
 \widehat f(\Lambda)=\sum_{v\in\Lambda\setminus\{0\}}f(v).
\]
For $p>0$, define a logarithmic fiber integral
\begin{align}
 V_f(p)
  =\sum_{\epsilon\in\{\pm1\}^6}\int_{\R^5}
  f\bigl(\epsilon_1e^{u_1},\ldots,\epsilon_5e^{u_5},
     \epsilon_6p e^{-u_1-\cdots-u_5}\bigr)
       \,\dd u_1\cdots\dd u_5 .                 \label{eq:fiber-integral}
\end{align}
Thus $V_f(p)$ integrates over the set on which the absolute product of
the six coordinates is $p$, with the Haar normalization used for $A$.

\begin{proposition}[Packet volume and vector unfolding]
\label{prop:packet-volume}
Every orbit in $\mathcal P_{K,\sigma}$ has $A$-volume $(s/2)R_K$,
and there are $h_K2^6/s$ distinct orbits.  In particular, the total
packet volume is $Q\kappa$.  For every nonnegative Borel function $f$,
\begin{equation}
 \int_X\widehat f\,\dd\mu_{K,\sigma}
   =\frac1{Q\kappa}\sum_{l\ge1}d_K(l)V_f(l/Q),
                                                        \label{eq:vector-unfolding}
\end{equation}
where both sides are allowed to be infinite.
\end{proposition}

\begin{proof}
The only roots of unity in a totally real field are $1$ and $-1$.
The kernel of the absolute logarithmic map on $U$ therefore has
order two.  Its image on $U^+$ has index $s/2$ in its image on $U$.
The unit theorem now gives covolume $(s/2)R_K$ for the positive-unit
logarithms in the coordinates $u_1,\ldots,u_5$.

We check that these are exactly the stabilizers and identifications
in the packet.  Suppose a real diagonal map sends $\sigma(I)$ onto
$\sigma(J)$.  Choose $0\ne d\in I$, and let its image be $\sigma(e)$
with $e\in J$.  All embeddings of $d$ are nonzero, so the diagonal
entries of the map are $\sigma_j(e/d)$.  It is multiplication by
$e/d\in K^\times$, and hence $J=(e/d)I$.  This observation also
applies to normalized lattices after their positive covolume factors
have been undone.  It follows that different ordinary ideal classes
give disjoint packet orbits.  For a fixed ideal, the diagonal
automorphisms are its units, because $(I:I)=\mathcal O_K$.

The stabilizer in $A$ is consequently the image of $U^+$.
Two sign translates belong to the same $A$-orbit exactly when their
signs differ by a unit signature.  The signature image has cardinality
$s$, giving $2^6/s$ orbits for each ordinary class and the stated orbit
volume.  The analytic class number formula, in this normalization, is
\[
 Q\kappa=2^5h_KR_K.
\]
It agrees with the product of the orbit number and orbit volume.

For completeness, let $\mathcal F$ be a fundamental domain for the
positive-unit logarithms in the traceless diagonal algebra.  Summing
over ordinary ideal-class representatives and all sign matrices
counts each orbit exactly $s$ times.  Thus for any nonnegative $F$ on
$X$,
\begin{equation}
 \int_XF\,\dd\mu_{K,\sigma}
 =\frac1{sQ\kappa}\sum_{[I]}\sum_w
   \int_{\mathcal F}F(\Lambda_{I,\sigma}w a(u))\,\dd u.
                                                        \label{eq:packet-integral}
\end{equation}
Here $\dd u=\dd u_1\cdots\dd u_5$ and the sixth coordinate is minus
their sum.

Apply this identity to $F=\widehat f$.  A nonzero $d\in I$ determines
the integral ideal
\[
 \mathfrak b=(d)I^{-1},
 \qquad
 \prod_{j=1}^6
 \left|\covol(\sigma(I))^{-1/6}\sigma_j(d)\right|
 =\frac{N\mathfrak b}{Q},
\]
because $\covol(\sigma(I))=QNI$.  The ideal $\mathfrak b$
determines the class $[I]=[\mathfrak b]^{-1}$, and every integral
$\mathfrak b$ occurs.  For that ideal, the possible $d$ form a single
$U$-orbit, or $s$ orbits under $U^+$.  Summing a positive-unit orbit
and integrating over $\mathcal F$ unfolds to the entire traceless
diagonal algebra.  Multiplication by a representative of $U/U^+$
translates the absolute logarithms and permutes the signs.  After the
sum over $w$, each of these $s$ orbits therefore contributes
$V_f(N\mathfrak b/Q)$.  Their number cancels the factor $s$ in
\eqref{eq:packet-integral}, proving \eqref{eq:vector-unfolding}.
All rearrangements are justified by nonnegativity.
\end{proof}

\subsection{Uniform counts of ideals}

The normalization by $\kappa$ in \eqref{eq:vector-unfolding} makes it
important to keep this same factor in upper bounds for ideal counts.
We obtain it from a zero-free neighborhood of $1$, followed by Shiu's
estimate for nonnegative multiplicative functions.

\begin{lemma}[Uniform ideal estimates]\label{lem:ideal-estimate}
For all sufficiently large discriminants $D$ in the present family,
\begin{align}
 \kappa&\gg (\log D)^{-1},                         \label{eq:residue-lower}\\
 \sum_{p\le D}\frac{d_K(p)}p
   &\le \log(\kappa\log D)+O(1),                 \label{eq:prime-sum}\\
 \sum_{x-y<l\le x}d_K(l)&\ll y\kappa
 \quad\left(D^{1/4}\le x\le D,
                    \ x^{1/2}\le y\le x\right).  \label{eq:ideal-interval}
\end{align}
In \eqref{eq:prime-sum}, $p$ ranges over rational primes.  The constants
and the discriminant threshold are uniform in $K$.
\end{lemma}

\begin{proof}
We first record the precise zero exclusion needed here.  Stark's
Lemma~3 shows that the rectangle
\[
 \Re s\ge1-\frac1{4\log D},\qquad
 |\Im s|\le\frac1{4\log D}
\]
contains at most one nontrivial zero of $\zeta_K$, and that such a
zero is real and simple.  His Lemma~8 applies also to nonnormal
fields: a real zero $\beta$ with
\[
 1-\frac1{4\cdot6!\log D}\le\beta<1
\]
would be a zero of the zeta function of a quadratic subfield of $K$
\cite[Lemmas~3 and~8]{Stark1974}.  Our field has no such subfield.
Consequently, for an absolute $c_0>0$,
\begin{equation}
 |1-\rho|\ge \frac{c_0}{\log D}
 \quad\text{for every nontrivial zero $\rho$ of $\zeta_K$}.
                                                        \label{eq:zero-exclusion}
\end{equation}
One can take $c_0=1/2880$.  This use of Stark's result does not require
$K/\Q$ to be normal.

Set $s_0=1+1/\log D$.  The logarithmic derivative of the completed
Dedekind zeta function gives, for real $1<s\le s_0$,
\begin{equation}
 S(s):=\sum_\rho\Re\frac1{s-\rho}
   =\frac1{s-1}+\frac12\log D
        +\frac{\zeta'_K(s)}{\zeta_K(s)}+O(1).
                                                        \label{eq:zeta-derivative}
\end{equation}
The zeros are counted with multiplicity.  The real-part sum converges,
and the gamma-factor contribution is bounded uniformly because the
degree and signature are fixed.  Each summand on the left is
nonnegative, and $\zeta'_K(s_0)/\zeta_K(s_0)\le0$ by its Dirichlet
series.  Hence $S(s_0)\ll\log D$.

The same estimate holds throughout $(1,s_0]$.  Indeed, writing
$a=1-\Re\rho\ge0$, $v=\Im\rho$, and $u=s-1$, the relevant term is
\[
 \frac{a+u}{(a+u)^2+v^2}.
\]
For $0<u\le1/\log D$, its numerator does not exceed its numerator
at $s_0$, while \eqref{eq:zero-exclusion} gives
$a^2+v^2\ge c_0^2/(\log D)^2$.  The denominator at $s_0$ is
therefore at most a constant multiple of the denominator at $s$.
Termwise comparison proves the asserted uniform estimate for $S(s)$.
Integrating \eqref{eq:zeta-derivative} after moving $1/(s-1)$ to the
other side now gives
\begin{equation}
 \zeta_K(s_0)\asymp\kappa\log D.                 \label{eq:zeta-residue-comparison}
\end{equation}
More explicitly, the logarithmic derivative of
$(s-1)\zeta_K(s)$ is $S(s)-\tfrac12\log D+O(1)$, whose integral
over this interval is $O(1)$, and its limit at $s=1$ is $\kappa$.
Since $\zeta_K(s_0)\ge1$, this proves \eqref{eq:residue-lower}.

The Euler product gives
\[
 \sum_{p\le D}\frac{d_K(p)}{p^{s_0}}
      \le\log\zeta_K(s_0).
\]
Also $d_K(p)\le6$ and
\[
 \sum_{p\le D}d_K(p)\left(\frac1p-\frac1{p^{s_0}}\right)
 \le\frac6{\log D}\sum_{p\le D}\frac{\log p}{p}\ll1,
\]
where the final estimate follows from Chebyshev's bound by partial
summation.  Equation~\eqref{eq:zeta-residue-comparison} now proves
\eqref{eq:prime-sum}.

To obtain the interval estimate, we use Shiu's theorem
\cite{Shiu1980}, in the formulation of
\cite[Theorem~1.1]{Pollack2020} with empty excluded residue sets.
For a nonnegative multiplicative function $f$, the hypotheses are
\[
 f(p^a)\le A_1^a,
 \qquad
 f(m)\le A_2(\epsilon)m^\epsilon
       \quad(\epsilon>0).
\]
For any fixed $0<\beta<1/2$, the conclusion is
\begin{equation}
 \sum_{x-y<m\le x}f(m)
 \ll\frac y{\log x}
     \exp\left(\sum_{p\le x}\frac{f(p)}p\right)
 \qquad(x^\beta<y\le x).                         \label{eq:shiu-form}
\end{equation}
Both the constant and the sufficiently-large threshold for $x$
depend only on $\beta,A_1$, and the function $A_2$.

The coefficients $d_K$ are multiplicative and are bounded by the
ordinary six-fold divisor function $d_6$.  In particular,
\[
 d_K(p^a)\le\binom{a+5}{5}\le6^a,
 \qquad d_K(m)\le d_6(m)\ll_\epsilon m^\epsilon.
\]
The hypotheses of \eqref{eq:shiu-form} thus hold with constants
independent of $K$.  We take, for example, $\beta=1/3$; the assumed
length $y\ge x^{1/2}$ lies in the required range.  Since $x\le D$,
\eqref{eq:prime-sum} bounds the exponential factor by
$O(\kappa\log D)$.  Finally $x\ge D^{1/4}$ implies
$\log D/\log x\le4$, proving \eqref{eq:ideal-interval}.
\end{proof}

\subsection{The cusp and balls in the space of lattices}

For a lattice $x\in X$, let $m(x)$ be the length of its shortest
nonzero vector in the sup norm on $\R^6$.  We use the sup norm only
to simplify the fiber integral; replacing it by any fixed norm
changes the constants.

\begin{proposition}[Uniform cusp bound]\label{prop:cusp}
There are $C>0$ and $c>0$ such that, for all sufficiently large $D$,
\[
 \mu_{K,\sigma}\{x:m(x)\le r\}\le Cr^c
       \qquad(0<r<1).
\]
Consequently the packet probabilities form a tight family as
$D\to\infty$, and every weak subsequential limit is an
$A$-invariant probability measure on $X$.
\end{proposition}

\begin{proof}
Take $f=\one_{[-r,r]^6}$ and put $X_r=Qr^6$.  For $0<p\le r^6$,
the change of variables $t_j=\log r-u_j$ in
\eqref{eq:fiber-integral} identifies the positive-sign integration
region with
\[
 t_j\ge0\quad(1\le j\le5),\qquad
 \sum_{j=1}^5t_j\le\log(r^6/p).
\]
Thus
\[
 V_f(p)=\frac{2^6}{5!}\bigl(\log(r^6/p)\bigr)^5
       \quad(0<p\le r^6),
 \qquad V_f(p)=0\quad(p>r^6).
\]
Since the occurrence of a short vector is bounded by its count,
\eqref{eq:vector-unfolding} gives
\begin{equation}
 \mu_{K,\sigma}\{m\le r\}
 \ll\frac1{Q\kappa}
       \sum_{l\le X_r}d_K(l)
                   \bigl(1+\log(X_r/l)\bigr)^5.          \label{eq:cusp-sum}
\end{equation}
The sum is empty if $X_r<1$.

Suppose first that $r\ge D^{-1/48}$, so $X_r\ge D^{3/8}$.
For $l\ge D^{1/4}$, divide the range of summation into intervals
$(2^{-j-1}X_r,2^{-j}X_r]$.  Whenever such an interval contains one
of these $l$, its upper endpoint is at least $D^{1/4}$ and at most
$Q$.  Lemma~\ref{lem:ideal-estimate}, with interval length half the
upper endpoint, bounds its ideal count by
$O(\kappa 2^{-j}X_r)$.  Its weight in \eqref{eq:cusp-sum} is
$O((j+1)^5)$.  The total contribution of these terms before division
by $Q\kappa$ is therefore at most
\[
 C\kappa X_r\sum_{j\ge0}2^{-j}(j+1)^5\ll\kappa X_r.
\]
For the remaining terms, the elementary divisor estimate
$\sum_{l\le Y}d_6(l)\ll Y(\log(2Y))^5$ gives an upper bound
$CD^{1/4}(\log D)^{10}$ before normalization.  In view of
\eqref{eq:residue-lower}, its normalized contribution is at most
\[
 CD^{-1/4}(\log D)^{11}\ll r^6;
\]
the last inequality follows from $r^6\ge D^{-1/8}$.
This proves an $O(r^6)$ bound in the present range.

Now let $r<D^{-1/48}$ and $X_r\ge1$.  Dyadic summation of the
same elementary divisor estimate yields
\[
 \sum_{l\le X_r}d_6(l)\bigl(1+\log(X_r/l)\bigr)^5
      \ll X_r(\log(2X_r))^5.
\]
Using $X_r\le Q$ and \eqref{eq:residue-lower},
\eqref{eq:cusp-sum} is therefore at most
$Cr^6(\log D)^6$.  But $\log D<48\log(1/r)$ in this range, and
$r^6(\log(1/r))^6\ll r^3$.  The claimed estimate follows, for
example with $c=3$, after increasing $C$.

Mahler's compactness criterion says that
$\{x:m(x)\ge r\}$ is compact for every $r>0$.  The bound just proved
therefore gives tightness.  Any weak limit has mass one, and
$A$-invariance passes to the limit because right translation by each
fixed element of $A$ is continuous and preserves compact supports.
\end{proof}

Let $B(r)$ denote an open matrix ball about the identity in $G$, for
$r$ small; all fixed smooth local metrics give equivalent estimates.

\begin{proposition}[Ordinary balls in a probability limit]
\label{prop:ball}
Suppose $D_i\to\infty$ and $\mu_i$ converges weakly to a probability
$\mu$.  For every compact $\Omega\subset X$, there are constants
$C_\Omega,r_\Omega>0$ such that
\[
 \mu(xB(r))\le C_\Omega r^6
       \qquad(x\in\Omega,\ 0<r<r_\Omega).
\]
\end{proposition}

\begin{proof}
At any lattice $x_0$, choose a nonzero lattice vector none of whose
coordinates is zero.  Such a vector exists because a full-rank lattice
cannot be contained in a finite union of proper linear subspaces.
In an injective local chart at $x_0$, continue this marked vector as
$v(x)$.  By reducing the chart, its coordinates remain bounded and
bounded away from zero.  A finite collection of these charts covers
$\Omega$, so these bounds and all local derivative bounds can be
taken uniform.

For $x$ in one of the charts and sufficiently small $r$, every lattice
in $xB(r)$ has a nonzero vector in the ordinary Euclidean ball
$B_{\R^6}(v(x),C_1r)$.  Let $f_{x,r}$ be the indicator of that ball.
Its support has absolute coordinate product in an interval $J_{x,r}$ of length
at most $C_2r$; these intervals all lie in a fixed compact subinterval
of $(0,\infty)$ when $r<r_\Omega$.  On this support the first five
logarithmic coordinates each vary by $O(r)$.  Consequently
\[
 V_{f_{x,r}}(p)\le C_3r^5\one_{J_{x,r}}(p).
\]
Fix $r>0$.  For all sufficiently large $i$, uniformly in the finitely
many charts and in $x\in\Omega$, the ideal norms in $Q_iJ_{x,r}$
can be covered by an interval with upper endpoint comparable to $Q_i$
and length at most $C_4Q_ir$.  We may choose that length comparable
to $Q_ir$, so it is at least the square root of its upper endpoint
for large $i$.  The upper endpoint lies between $D_i^{1/4}$ and
$D_i$.  Lemma~\ref{lem:ideal-estimate} and
\eqref{eq:vector-unfolding} therefore give
\[
 \mu_i(xB(r))
 \le\int_X\widehat f_{x,r}\,\dd\mu_i
 \le\frac{C_3r^5}{Q_i\kappa_i}
       \sum_{l\in Q_iJ_{x,r}}d_{K_i}(l)
 \le C_5r^6.
\]
The constants are independent of $x$ and $r$ in the stated range;
the threshold for $i$ may depend on $r$.  Since $xB(r)$ is open,
the open-set inequality for weak convergence gives the same bound
for $\mu(xB(r))$.  This proves the assertion for each $x$ and $r$
with the uniform constants supplied by the finite cover.
\end{proof}

\subsection{Trace slices of the codifferent}

The vector estimate has now supplied tightness and the ordinary-ball
bound.  To control diagonal tubes, we need to count a quantity which
is nearly preserved by block-diagonal matrices.  For an ordinary
vector $v$ and a dual vector $y$, that quantity is a partial sum of
the products $v_jy_j$.  The following lattice estimate makes such
partial sums effective arithmetic coordinates.

Let $\mathfrak D_K$ be the different of $K$, so
\[
 \mathfrak D_K^{-1}
 =\{z\in K:\Tr_{K/\Q}(z\mathcal O_K)\subset\Z\}.
\]
Give $\R^6$ its Euclidean inner product, let
$\mathbf e=(1,\ldots,1)$, and put
\[
 H=\mathbf e^\perp,
 \qquad \mathcal L_K=\sigma(\mathfrak D_K^{-1})\cap H.
\]

\begin{lemma}[Small cells in a trace slice]\label{lem:trace-slice}
The lattice $\mathcal L_K$ has Euclidean covolume $\sqrt6/Q$ in
$H$ and admits a fundamental parallelotope of diameter
$O(D^{-1/30})$.  Let $S$ be a nonempty proper subset of
$\{1,\ldots,6\}$.  For fixed $M>0$, uniformly in $m\in\Z$ and
intervals $J\subset\R$ of length $\delta>0$,
\begin{equation}
 \#\left\{z\in\mathfrak D_K^{-1}:
     \Tr_{K/\Q}z=m,
     \ |\sigma_j(z)|\le M\ (1\le j\le6),
     \ \sum_{j\in S}\sigma_j(z)\in J\right\}
 \ll_M Q\bigl(\delta+D^{-1/30}\bigr).
                                                        \label{eq:trace-slice-count}
\end{equation}
\end{lemma}

\begin{proof}
Since $1$ is primitive in the free abelian group $\mathcal O_K$,
choose an integral basis $\alpha_1=1,\alpha_2,\ldots,\alpha_6$.
Let $\beta_1,\ldots,\beta_6$ be its trace-dual basis, so
$\Tr_{K/\Q}(\alpha_j\beta_k)=\delta_{jk}$.  By the definition of
the codifferent, these $\beta_k$ form a basis of
$\mathfrak D_K^{-1}$.  In particular,
$\Tr_{K/\Q}\beta_k=\delta_{1k}$, which proves that the trace map
onto $\Z$ is surjective and that
$\sigma(\beta_2),\ldots,\sigma(\beta_6)$ form a basis of
$\mathcal L_K$.  For $2\le j,k\le6$, orthogonal projection gives
\[
 \bigl\langle\operatorname{proj}_H\sigma(\alpha_j),
                    \sigma(\beta_k)\bigr\rangle=\delta_{jk}.
\]
These projected vectors are therefore the dual basis inside $H$.
Since $\operatorname{proj}_H\sigma(\alpha_1)=0$, we obtain
\begin{equation}
 \mathcal L_K^*=\operatorname{proj}_H\sigma(\mathcal O_K).
                                                        \label{eq:projected-dual}
\end{equation}

Take a nonzero vector in the right side of
\eqref{eq:projected-dual}, arising from $\alpha\in\mathcal O_K$.
The vector is nonzero precisely when $\alpha\notin\Q$.
Primitivity of the field implies $\Q(\alpha)=K$, so the order
$\Z[\alpha]$ has discriminant at least $D$.  Differences of
embedding coordinates are unchanged by orthogonal projection to
$H$.  It follows that
\[
 D\le |\Disc(\Z[\alpha])|
   =\prod_{j<k}|\sigma_j(\alpha)-\sigma_k(\alpha)|^2
   \le C\|\operatorname{proj}_H\sigma(\alpha)\|^{30}.
\]
Thus the first successive minimum of $\mathcal L_K^*$ is at least
$cD^{1/30}$.  Transference of successive minima in dimension five
gives
\[
 \lambda_5(\mathcal L_K)
       \ll\lambda_1(\mathcal L_K^*)^{-1}
       \ll D^{-1/30}.
\]
A fixed-dimensional lattice basis bound then provides a basis for
$\mathcal L_K$ all of whose lengths are $O(D^{-1/30})$; these
standard geometry-of-numbers facts may be found in \cite{Cassels1959}.
The associated fundamental parallelotope has the asserted diameter.

The full codifferent lattice has covolume $Q^{-1}$.  Its trace
planes have spacing $1/\sqrt6$, since the trace takes every integer
value and has normal vector $\mathbf e$.  Decomposing a basis into
five vectors in $H$ and a vector of trace one therefore gives
\[
 Q^{-1}=\frac1{\sqrt6}\covol_H(\mathcal L_K).
\]

The trace-$m$ points, if any, form a translate of $\mathcal L_K$.
Tile that affine hyperplane by translates of the fundamental
parallelotope, assigning boundaries measurably.  A cell attached to
a point counted in \eqref{eq:trace-slice-count} lies in the
$O(D^{-1/30})$ enlargement of its bounded slice.  The linear
functional $z\mapsto\sum_{j\in S}z_j$ is nonconstant on $H$ because
$S$ is nonempty and proper.  The five-dimensional volume of the
enlarged slice is consequently
$O_M(\delta+D^{-1/30})$, uniformly in its location and in $m$.
Dividing by the cell covolume $\sqrt6/Q$ proves the counting bound.
\end{proof}

\subsection{Exponential decay on diagonal tubes}

For $D_i$ sufficiently large, define
\begin{equation}
 T_{-,i}=\frac{\log D_i}{(\log\log D_i)^{2/3}},\qquad
 T_{+,i}=\frac{\log D_i}{(\log\log D_i)^{1/3}}.
                                                        \label{eq:time-range}
\end{equation}
Both tend to infinity, their ratio tends to infinity, and
$T_{+,i}=o(\log D_i)$.  The last property ensures that the cells in
Lemma~\ref{lem:trace-slice} are much smaller than the slices required
at these times.

\begin{proposition}[Two-sided tube estimate]\label{prop:bowen}
Let $S$ be a nonempty proper subset of $\{1,\ldots,6\}$, and let
$b=(b_1,\ldots,b_6)\in L$ satisfy
\[
 |b_j-b_k|\ge1\qquad(j\in S,\ k\notin S).
\]
For every compact $\Omega\subset X$, there is an identity
neighborhood $O\subset G$ such that, for all sufficiently large
$i$, all $x\in\Omega$, and all $T_{-,i}\le t\le T_{+,i}$,
\begin{equation}
 \mu_i\bigl(x(a(tb)Oa(-tb)\cap a(-tb)Oa(tb))\bigr)
       \le e^{-t/3}.                              \label{eq:bowen-bound}
\end{equation}
The neighborhood and discriminant threshold may be chosen
simultaneously for any fixed finite collection of pairs $(S,b)$.
\end{proposition}

\begin{proof}
For a lattice $\Lambda$, its Euclidean dual is
$\Lambda^*=\{y:\langle v,y\rangle\in\Z\text{ for all }v\in\Lambda\}$.
At each point of $\Omega$, choose nonzero marked vectors in the
lattice and its dual, both avoiding all coordinate hyperplanes.
Continue the pair in a local lattice chart.  Finitely many such
charts cover $\Omega$, and the marked vectors in these charts lie
in fixed bounded boxes separated from the coordinate hyperplanes.
Their inner product is an integer, constant in each chart.  Indeed,
under right multiplication by $h$, the pair transforms as
$v\mapsto vh$, $y\mapsto yh^{-\mathsf T}$.

Choose $O$ sufficiently small in matrix coordinates.  If
\[
 h\in a(tb)Oa(-tb)\cap a(-tb)Oa(tb),
\]
then $h$ is uniformly close to the identity, and each entry crossing
the cut $S\sqcup S^c$ is $O(e^{-t})$.  To verify the latter bound,
the two membership conditions respectively bound its $(j,k)$ entry
by a constant times $e^{t(b_j-b_k)}$ and $e^{-t(b_j-b_k)}$;
their minimum is $O(e^{-t|b_j-b_k|})$.  Reducing $O$ also makes
$h^{-1}$ uniformly bounded.  Let $P_S$ be the diagonal projection
onto the coordinates in $S$.  Then
\[
 hP_Sh^{-1}-P_S=[h,P_S]h^{-1}=O(e^{-t}).
\]
For the marked pair $v_0,y_0$ at the center $x$, put
$v=v_0h$ and $y=y_0h^{-\mathsf T}$.  These vectors remain in fixed
bounded boxes avoiding the coordinate hyperplanes, and
\begin{align*}
 \sum_{j=1}^6v_jy_j&=\langle v_0,y_0\rangle=:m\in\Z,\\
 \sum_{j\in S}v_jy_j
   &=v_0hP_Sh^{-1}y_0^{\mathsf T}
     =v_0P_Sy_0^{\mathsf T}+O(e^{-t}).
\end{align*}
All constants are uniform over the finite cover of $\Omega$.
Thus membership in the tube forces the existence of an ordinary
and dual vector pair satisfying these boundedness, trace, and
partial-sum conditions.  The short partial-sum interval may depend
on $x$.

We count such pairs in the packet.  For a representative ideal $I$
and a packet lattice $c\sigma(I)wa$, where
$c=\covol(\sigma(I))^{-1/6}$, its dual is
\[
 c^{-1}\sigma(\mathfrak D_K^{-1}I^{-1})wa^{-1}.
\]
Consequently a nonzero pair comes from
\[
 d\in I\setminus\{0\},\qquad
 \ell\in\mathfrak D_K^{-1}I^{-1}\setminus\{0\}.
\]
The scalar, sign, and diagonal factors cancel in their
coordinatewise product:
\begin{equation}
 v_jy_j=\sigma_j(z),\qquad z=d\ell\in\mathfrak D_K^{-1}.
                                                        \label{eq:pair-product}
\end{equation}
The allowed $z$ therefore belong to a fixed bounded embedding box,
have trace $m$, and have their $S$-sum in an interval of length
$O(e^{-t})$.

For a fixed nonzero allowed $z$, the ordinary vector's ideal
$\mathfrak b=(d)I^{-1}$ must divide $(z)\mathfrak D_K$.  In fact
\[
 \mathfrak b\bigl((\ell)I\mathfrak D_K\bigr)
       =(z)\mathfrak D_K,
\]
and both factors on the left are integral.  Conversely, this
divisibility ensures $z/d\in\mathfrak D_K^{-1}I^{-1}$ whenever
$(d)I^{-1}=\mathfrak b$.  For a fixed $d$ and $z$, there is exactly
one possible $\ell$, namely $z/d$.

Let $f$ be the indicator of a fixed bounded box, or of a fixed finite
union of boxes, containing all possible ordinary vectors $v$ and
separated from the coordinate hyperplanes.  Its logarithmic fiber
integrals are bounded uniformly in $p>0$, because the first five
absolute logarithms range over a fixed bounded set.  Repeating
the unfolding of Proposition~\ref{prop:packet-volume}, with the
additional restriction $\mathfrak b\mid(z)\mathfrak D_K$, gives
\begin{equation}
 \int_X\#\{\text{admissible pairs with product }z\}
           \,\dd\mu_{K,\sigma}
 \le\frac1{Q\kappa}
       \sum_{\mathfrak b\mid(z)\mathfrak D_K}
          V_f(N\mathfrak b/Q)
 \ll\frac1{Q\kappa}
          \#\{\mathfrak b:\mathfrak b\mid(z)\mathfrak D_K\}.
                                                        \label{eq:pair-unfolding}
\end{equation}
Here we have dropped the restrictions on the dual vector after
fixing $z$, which can only increase the count.  The divisibility
restriction is unchanged by multiplication of $d$ by a unit, so
the same unfolding and sign multiplicities apply.

The embedding bound on $z$ implies
$N((z)\mathfrak D_K)=D|N_{K/\Q}z|\ll D$.
For an integral ideal of norm $N$, its number of ideal divisors is
at most $d(N)^6$, where $d$ is the ordinary divisor function: over
each rational prime there are at most six prime ideals, and every
one of their exponents is at most the exponent of that prime in
$N$.  The classical divisor bound consequently gives, uniformly
for our allowed $z$,
\begin{equation}
 \#\{\mathfrak b:\mathfrak b\mid(z)\mathfrak D_K\}
       \le\exp\left(C\frac{\log D}{\log\log D}\right).
                                                        \label{eq:ideal-divisors}
\end{equation}

Lemma~\ref{lem:trace-slice} bounds the number of allowed products
by $O(Q(e^{-t}+D^{-1/30}))$.  Uniformly for $t\le T_{+,i}$,
\[
 D_i^{-1/30}=o(e^{-t}),
\]
because $T_{+,i}/\log D_i\to0$.  The product count is therefore
$O(Q_i e^{-t})$.  The event defining the tube is bounded by the
number of admissible pairs, so summing \eqref{eq:pair-unfolding}
over these products and using \eqref{eq:ideal-divisors} gives
\[
 \mu_i\bigl(x(a(tb)Oa(-tb)\cap a(-tb)Oa(tb))\bigr)
 \ll\kappa_i^{-1}
      \exp\left(-t+C\frac{\log D_i}{\log\log D_i}\right).
\]
Finally, \eqref{eq:residue-lower} and $t\ge T_{-,i}$ show that
the logarithm of the factor multiplying $e^{-t}$ is $o(t)$,
uniformly throughout the time range.  Enlarging the discriminant
threshold makes the last display at most $e^{-t/3}$.
The same choices work for finitely many cuts and diagonal elements
by taking the smallest of their identity neighborhoods and the
largest of their thresholds.
\end{proof}
